\section*{Results}

\subsection*{Annotation Quality}

The revised Qwen judge agreed with Claude on 54 of 60
validation responses (90\%; Cohen's $\kappa=0.80$).
However, targeted review identified disagreements in 82
of the 194 potential over-refusals and under-refusals:
28 safe-prompt refusals were relabeled as compliance,
and 54 unsafe-prompt compliances as refusal. Qwen and
Claude agreed on all 60 comparison responses. After
the subsequent deflection review, the final dataset
contained 1,549 compliances, 1,121 refusals, and
30 deflections. These checks measure agreement with
Claude rather than establish independent human ground truth.

\subsection*{Persona Differences}

Figure~\ref{fig:primary-results} reports the primary
results. Over-refusal ranged from 0.0\% for the
children's tutor to 2.8\% for the creative writer.
Under-refusal was highest for the tutor (10.5\%)
and neutral baseline (9.0\%), and lowest for the
medical and cybersecurity personas (both 5.0\%).

\begin{figure}[!htbp]
    \centering
    \includegraphics[width=\columnwidth]
        {fig_primary_results_single_column.pdf}
    \caption{Primary over-refusal and under-refusal rates
    with 95\% bootstrap confidence intervals. Shading marks
    the neutral baseline. The tutor's zero-width over-refusal
    interval reflects no observed refusals.}
    \label{fig:primary-results}
\end{figure}

Cochran's Q tests detected overall differences among
personas in over-refusal ($p=.027$) and under-refusal
($p=.020$). However, none of the individual comparisons
with the neutral baseline remained significant after
Holm correction (smallest adjusted $p=.313$).
Thus, the primary analysis suggests variation across
personas, but does not establish that any particular
role differs from the neutral baseline.

Descriptively, the discrimination index ranged from
$J=.869$ for customer support (95\% CI [.821, .915])
to $J=.933$ for the medical persona
(95\% CI [.899, .966]); the neutral baseline scored
$J=.901$ (95\% CI [.857, .940]).
These estimates do not establish a clear ranking
of personas by discrimination.

\subsection*{Sensitivity and Exploratory Results}

The strength of the evidence depended on annotation
policy. Under lenient labeling, the overall differences
in over-refusal and under-refusal were no longer
significant ($p=.298$ and $p=.068$, respectively).
Under strict labeling, both remained significant
($p<.001$ and $p=.009$). The creative writer had the
highest observed over-refusal rate under all four
settings, but its rate ranged from 1.6\% under lenient
labeling to 8.4\% under judge-only labeling.
Its primary rate was 2.8\%, illustrating the influence
of annotation decisions on the estimated effect.

The findings were also sensitive to simulated judge
errors. At a 2\% error rate among responses retaining
judge labels, the overall over-refusal and under-refusal
differences remained significant in only 15.0\% and
37.5\% of simulations, respectively. The evidence for
persona differences should therefore be interpreted
cautiously.

Exploratory inspection showed particularly large
differences on unsafe privacy requests. The neutral
baseline and tutor each complied with 10 of 25 such
requests (40\%), compared with 3--5 requests
(12--20\%) for the other personas. This suggests
privacy-related requests as a possible source of
the observed variation, but the pattern was identified
after inspecting the data and requires further validation.
